\section{椭圆的几何与极坐标}
\subsection{长短轴、焦点和面积}

椭圆的直角坐标方程为 $x^2/a^2+y^2/b^2=1$，其中 $a\ge b>0$。两焦点距中心均为 $c=\sqrt{a^2-b^2}$，离心率 $e=c/a$。把半径 $a$ 的圆沿一个方向压缩为原来的 $b/a$，得到该椭圆；相同投影同时把任意区域的面积乘以 $b/a$，所以整椭圆面积为 $\pi ab$。

取右焦点为引力中心 $F$，近日点 $P$ 在右端、远日点 $A$ 在左端。距离为 $r_p=a(1-e)$、$r_a=a(1+e)$，故两者之和是 $2a$。半长轴 $a$ 是两个极端距离的平均数，并不是任意时刻到焦点距离的平均值，也不是轨道上每处的曲率半径。

\typeout{ASTROFIG:fig:at2:ellipse}\fig{fig:at2:ellipse}{焦点、中心与近远地点分属不同几何位置}{\begin{tikzpicture}[scale=.85]
\draw[world](-1.5,0)ellipse(3 and 2.598076);\fill(0,0)circle(2pt)node[below right]{$F$};\fill(1.5,0)circle(1.5pt)node[right]{$P$};\fill(-4.5,0)circle(1.5pt)node[left]{$A$};\draw[guide](-4.5,0)--(1.5,0);\node[below]at(-1.5,0){$O$};\draw[guide](-1.5,0)--(-1.5,2.598);\node[left]at(-1.5,1.3){$b$};\draw[dim](-4.5,-.8)--(1.5,-.8)node[midway,below]{$2a$};
\end{tikzpicture}}
图\ref{fig:at2:ellipse}的半短轴从 $O$ 量起，引力距离从 $F$ 量起。短轴端点到焦点的距离恰为 $a$，这条几何关系在速度和时间问题中很有用。\subsection{焦点为极点时的轨道方程}

以 $FP$ 为极轴、逆时针极角为 $f$，椭圆可写成
\[r=\frac{p}{1+e\cos f},\qquad p=a(1-e^2)=\frac{b^2}{a}.\]
本讲用 $p$ 表示半通径，即 $f=\pi/2$ 时的极径。它不等于焦点到准线的距离；准线在焦点右侧的距离为 $p/e$，只有把定义中的比例 $e$ 一起乘回去，才能得到分子 $p$。当 $e=0$ 时轨道为圆，极坐标式仍成立，准线则移至无穷远。

还可用偏近点角 $E$ 参数化：相对于椭圆中心 $x_O=a\cos E,y=b\sin E$；相对于右焦点 $x_F=a(\cos E-e),y=b\sin E$。从勾股关系可得 $r=a(1-e\cos E)$。$E$ 是辅助圆上的角，不是焦点处看到的真近点角 $f$。两个角虽在近、远地点相同，在其他位置一般不同。


\subsection{从位置参数逐步求出时间参数}

  以右焦点为原点，令 $x=a(\cos E-e),y=b\sin E$。计算距离平方时，利用 $b^2=a^2(1-e^2)$，
  \begin{align*}
  r^2&=a^2(\cos E-e)^2+a^2(1-e^2)\sin^2E\\
  &=a^2(1-2e\cos E+e^2\cos^2E)=a^2(1-e\cos E)^2.
  \end{align*}
  椭圆 $e<1$ 保证括号始终为正，因此 $r=a(1-e\cos E)$。这里取正根有几何依据，不是随意丢掉绝对值。

  焦点扫面积的微分可由两个位置矢量的叉积给出：$\dd S=(x\dd y-y\dd x)/2$。将 $\dd x=-a\sin E\dd E$、$\dd y=b\cos E\dd E$ 代入，得到
  \[\dd S=\frac{ab}{2}(1-e\cos E)\dd E.
  \]
  它等于面积速度 $h/2$ 乘 $\dd t$，所以 $\dd t=(ab/h)(1-e\cos E)\dd E$，再积分才得到开普勒时间式。这一步说明时间权重并非人为选取，而是焦点处的面积守恒自然给出的。
  
\typeout{ASTROFIG:fig:at2:parameter}\fig{fig:at2:parameter}{同一轨道点的中心参数与焦点径矢相互对应}{\begin{tikzpicture}[scale=.85]
\draw[world](-1.5,0)ellipse(3 and 2.598076);\fill(0,0)circle(2pt)node[below right]{$F$};\fill(1.5,0)circle(1.5pt)node[right]{$P$};\fill(-4.5,0)circle(1.5pt)node[left]{$A$};\draw[guide](-4.5,0)--(1.5,0);\node[below]at(-1.5,0){$O$};\fill(-1.5,2.598)circle(2pt)node[above]{$B$};\draw[guide](-1.5,0)--(-1.5,2.598);\draw[force](0,0)--(-1.5,2.598);\node at(-3,1.2){$E$ 从中心计};\node at(.7,1.5){$f$ 从焦点计};
\end{tikzpicture}}

  图\ref{fig:at2:parameter}把两种角的参考点分开。把 $x/r,y/r$ 写成 $\cos f,\sin f$，便得到
  \[\cos f=\frac{\cos E-e}{1-e\cos E},\qquad
  \sin f=\frac{\sqrt{1-e^2}\sin E}{1-e\cos E}.\]
  两式给出同一个点的完整方向信息。只取第一式反余弦会把上半轨道和下半轨道混淆；处理完整一周时，可以连续跟踪角度，也可以使用同时输入正弦、余弦的带象限反正切。

  速度同样可以先对 $E$ 求导，再除以 $\dd t/\dd E$，得到 $v_x=-an\sin E/(1-e\cos E)$、$v_y=bn\cos E/(1-e\cos E)$。当 $E=0$ 时速度向上，当 $E=\pi$ 时速度向下；在两个短轴端，竖直分量为零。检查这些特殊位置，能及时发现角度、焦点方向或时间分支选错。

  给定时间反求位置时，先解 $E-e\sin E=M$，再计算 $x,y$，最后按需要转成极角。把这个过程反过来，直接用 $f=nt$ 画椭圆上的等时间点，会错误地把极角等分。正确图中，近日点附近相邻等时间位置相隔较远，远日点附近相隔较近；这既是面积定律的图像，也是数值程序输出的必要检查。

  若轨道参数来自测量，速度、半长轴与周期的误差也会相互传递。例如 $T\propto a^{3/2}$，小误差有 $\Delta T/T\simeq3\Delta a/(2a)$。使用这一关系前，中心引力参数须固定；若 $GM$ 也有误差，应补上 $-\Delta\mu/(2\mu)$，不能只保留半长轴的一项。
  \section{角动量、能量与轨道形状}
\subsection{中心力首先给出面积速度}

记引力参数 $\mu=GM$，小卫星质量相对中心天体可忽略。引力沿径向，关于焦点的力矩为零，所以单位质量角动量 $h=r^2\dot f$ 为常量。很短时间内扫过的扇形面积为 $r^2\dd f/2$，因此面积速度为 $h/2$。这就是开普勒第二定律，它不要求速度大小恒定。

将速度分成径向 $\dot r$ 和横向 $r\dot f=h/r$，单位质量机械能为
\[\varepsilon=\frac12\dot r^2+\frac{h^2}{2r^2}-\frac\mu r.
\]
后两项可看成径向运动的有效势。初态给定 $r,v_r,v_t$，应先算 $h=rv_t$ 与 $\varepsilon=(v_r^2+v_t^2)/2-\mu/r$；只给总速率能确定能量，却不能确定角动量或最近距离。

\subsection{从动力学求半通径和离心率}

令 $u=1/r$，以 $f$ 为自变量。由 $\dot f=hu^2$，链式求导得到 $\dot r=-hu'$、$\ddot r=-h^2u^2u''$。径向牛顿方程 $\ddot r-r\dot f^2=-\mu/r^2$ 化为
\[u''+u=\frac\mu{h^2}.\]
把近地点选为 $f=0$，解为 $u=(\mu/h^2)(1+e\cos f)$，所以 $p=h^2/\mu$。再代入能量式，$u'=-(\mu e/h^2)\sin f$，有
\[\varepsilon=\frac{\mu^2}{2h^2}(e^2-1),\qquad
 e=\sqrt{1+\frac{2\varepsilon h^2}{\mu^2}}.
\]
这一步使几何参数与初始状态相连。$\varepsilon<0$ 给椭圆，零给抛物线，正给双曲线；但 $h=0$ 是退化的径向运动，必须直接用径向能量积分，不能把极角当作仍在变化的自变量。

\subsection{椭圆速度与两体推广}

椭圆满足 $p=a(1-e^2)$，上式给 $\varepsilon=-\mu/(2a)$，于是任意位置的速度满足
\[v^2=\mu\left(\frac2r-\frac1a\right).\]
近远地点的速度垂直径向，故 $r_pv_p=r_av_a=h$。在一般位置，守恒量是 $rv_t$，不能把总速率 $v$ 代替 $v_t$。

如果两颗质量都不可忽略，先在质心系求相对轨道，把 $\mu$ 换为 $G(m_1+m_2)$，半长轴 $a$ 指两体间相对距离轨道的半长轴。各颗星的位置都是相对位置乘一个固定比例，因此它们绕共同质心作相似椭圆，而不是绕对方的瞬时位置作固定焦点椭圆。


\section{运行时间由扫过面积决定}
\subsection{辅助圆扇形减去焦点三角形}

从近日点到偏近点角 $E$，中心扇形的投影面积为 $abE/2$。焦点与中心相差 $ae$，还要减去焦点三角形 $aeb\sin E/2$，所以扫过面积为 $ab(E-e\sin E)/2$。结合面积速度 $h/2$，得到
\[t-t_p=\frac{ab}{h}(E-e\sin E)=\frac1n(E-e\sin E),\quad n=\sqrt{\mu/a^3}.
\]
全周 $E=2\pi$，得到 $T=2\pi/n$。对于跨越多个象限的轨道，必须沿实际运动连续延伸角度，不能把反三角函数返回的主值机械代入。半周一定用时 $T/2$，半条弧长却未必用时 $T/4$。

\typeout{ASTROFIG:fig:at2:area}\fig{fig:at2:area}{辅助圆参数和焦点扫面积对应不同角}{\begin{tikzpicture}[scale=.85]
\draw[world](-1.5,0)ellipse(3 and 2.598076);\fill(0,0)circle(2pt)node[below right]{$F$};\fill(1.5,0)circle(1.5pt)node[right]{$P$};\fill(-4.5,0)circle(1.5pt)node[left]{$A$};\draw[guide](-4.5,0)--(1.5,0);\node[below]at(-1.5,0){$O$};\fill[blue!8](0,0)--(1.5,0)arc(0:90:3 and 2.598)--cycle;\draw[world](-1.5,0)--(-1.5,2.598);\fill(-1.5,2.598)circle(2pt)node[above]{$B$};\draw[guide](0,0)--(-1.5,2.598);
\end{tikzpicture}}
图\ref{fig:at2:area}的 $B$ 是短轴端，辅助角 $E=\pi/2$。焦点右移使近日点到 $B$ 的扫面积小于四分之一椭圆面积，因此用时小于 $T/4$。\subsection{从给定位置转换成时间}

已知半径 $r$ 时先由 $\cos E=(1-r/a)/e$ 求候选角；出近地点阶段取 $\sin E>0$，回程阶段取 $\sin E<0$。已知真近点角时可用
\[\cos E=\frac{e+\cos f}{1+e\cos f},\qquad
\sin E=\frac{\sqrt{1-e^2}\sin f}{1+e\cos f}.
\]
两式一起确定象限。若反过来已知时刻，应求 $E-e\sin E=n(t-t_p)$。椭圆 $e<1$ 时其左侧导数 $1-e\cos E>0$，每个连续周期内解唯一；可以用二分或牛顿法求数值，并以残差检查时间误差。

面积法适合椭圆弧段时间，能量积分适合径向上升与相落时间。两种方法都从同一动力学出发，但积分变量和几何边界不同；不能因为看到一个“最高点”，就把径向飞行当成半个非退化椭圆周期。


\section{速度方向、曲率与近似}
\subsection{曲率半径由法向加速度确定}

运动轨迹的法向加速度为 $v^2/\mathcal R$。引力一般不垂直速度，只有其法向分量提供这部分加速度。若速度与径向夹角为 $\psi$，则 $a_n=(\mu/r^2)|\sin\psi|$，所以 $\mathcal R=v^2r^2/(\mu|\sin\psi|)$。

在椭圆长轴端，速度恰垂直径向，结合活力公式得到 $\mathcal R=b^2/a$。短轴端切线水平，引力方向倾斜，法向分量为 $\mu b/a^3$，而 $v^2=\mu/a$，所以 $\mathcal R=a^2/b$。两个位置即使到焦点距离相近，曲率也可能很不一样。

接近圆轨道时 $e$ 为小参数，可比较 $r_p=a(1-e)$ 与 $r_a=a(1+e)$。若只求一阶变化，周期在固定 $a$ 时不变；若固定的是一个端点半径，$a$ 本身随 $e$ 改变，周期也改变。近似计算必须说明保持什么不变。


\section{轨道量的平均与速度图}
\subsection{时间平均必须按实际停留时间加权}

  卫星在远地点附近运行慢，不能把沿几何图形等分得到的算术平均当作时间平均。偏近点角给 $\dd t=(1-e\cos E)\dd E/n$，所以任意一周期平均为
  \[\langle A\rangle_t=\frac1{2\pi}\int_0^{2\pi}A(E)(1-e\cos E)\dd E.\]
  对 $A=1/r$，$r=a(1-e\cos E)$ 恰好抵消权重，得到 $\langle1/r\rangle=1/a$。再由活力式得 $\langle v^2\rangle=\mu/a$，故平均动能是 $\mu m/(2a)$，平均势能是 $-\mu m/a$。这不是说每一时刻动能都等于负势能的一半；只有圆轨道才逐点成立。

  对 $A=r$，权重不会抵消，而产生平方项。$\cos E$ 一周积分零、$\cos^2E$ 一周平均为一半，故 $\langle r\rangle=a(1+e^2/2)>a$。远处停留更久，使时间平均距离大于近远距离的简单平均 $a$。
  
\subsection{速度矢量的端点也有几何轨迹}

  用轨道极角 $f$ 写径向、横向分量，$v_r=(\mu e/h)\sin f$、$v_t=(\mu/h)(1+e\cos f)$。再投影到固定的近日点坐标轴，得到 $v_x=-(\mu/h)\sin f$、$v_y=(\mu/h)(e+\cos f)$。因此速度矢量端点在速度平面上落在一个圆上，圆心 $(0,\mu e/h)$、半径 $\mu/h$。

  这个圆位于速度平面，不是卫星的空间轨道。圆心偏移解释近远地点速度不相等，也给速度方向的极值提供了方便。建立图时应明确横纵轴的单位是速度，不能把椭圆空间图中的长度尺度直接搬进速度图。
  
\subsection{近圆轨道的时间与角度近似}

  对 $e\ll1$，开普勒方程 $E-e\sin E=M$ 的首阶解是 $E=M+e\sin M+O(e^2)$。由两种角的转换关系，$f=E+e\sin E+O(e^2)$，故 $f=M+2e\sin M+O(e^2)$，而 $r=a(1-e\cos M)+O(ae^2)$。

  平均近点角 $M$ 随时间匀速增加，真近点角 $f$ 则有一阶周期修正。在固定时刻比较不同轨道位置时，这个区别直接影响方向误差。将近圆轨道当作等角速度圆运动可以给零阶近似，但若还要求保留一阶偏差，就必须把时间参数的修正一起保留。
  

